\documentclass[11pt]{article}
\usepackage[utf8]{inputenc}
\usepackage[T1]{fontenc}
\usepackage{amsmath}
\usepackage{amssymb}
\usepackage{multicol}
\usepackage{geometry}
\usepackage{tikz}
\usepackage{enumitem}
\usepackage{xcolor}
\usepackage{titlesec}

\geometry{a4paper, left=1cm, right=1cm, top=0.5cm, bottom=1.2cm}
\setlength{\columnseprule}{0.4pt}

\titleformat{\section}{\normalfont\Large\bfseries\color{blue}}{}{0em}{}
\titleformat{\subsection}{\normalfont\large\bfseries\color{red}}{}{0em}{}

\title{\textcolor{red}{Medidas Estatísticas}}
\author{Professor: Jefferson}
\date{}

\begin{document}

\maketitle
\vspace{-1cm}

\begin{multicols}{2}

\section*{Exercícios com Resolução Comentada}

\begin{enumerate}[leftmargin=*]

% QUESTÃO 1
\item A média aritmética do conjunto $(x, 6, 10)$ é 8. Determine o valor de $x$.

\textbf{Resolução:}

A média é dada por:
\[
\bar{x} = \frac{x + 6 + 10}{3}
\]

Como a média vale 8:
\[
\frac{x + 16}{3} = 8
\Rightarrow x + 16 = 24
\Rightarrow x = 8
\]

\textbf{Resposta:} $x = 8$.

% QUESTÃO 2
\item A média dos números $(2x, 4, 6)$ é 10. Calcule $x$.

\textbf{Resolução:}

\[
\frac{2x + 4 + 6}{3} = 10
\Rightarrow \frac{2x + 10}{3} = 10
\Rightarrow 2x + 10 = 30
\Rightarrow 2x = 20
\Rightarrow x = 10
\]

\textbf{Resposta:} $x = 10$.

% QUESTÃO 3
\item O conjunto $(x, 5, 7, 9)$ possui média igual a 7. Determine $x$.

\textbf{Resolução:}

\[
\frac{x + 5 + 7 + 9}{4} = 7
\Rightarrow \frac{x + 21}{4} = 7
\Rightarrow x + 21 = 28
\Rightarrow x = 7
\]

\textbf{Resposta:} $x = 7$.

% QUESTÃO 4
\item A mediana do conjunto $(2, x, 8, 10, 12)$ é 8. Determine $x$.

\textbf{Resolução:}

Como o conjunto possui 5 elementos, a mediana é o termo central.
Para que o termo central seja 8, o valor de $x$ deve ser menor ou igual a 8.

O único valor que mantém a ordem correta é:
\[
x = 8
\]

\textbf{Resposta:} $x = 8$.

% QUESTÃO 5
\item A moda do conjunto $(4, x, 6, x, 8)$ é 6. Determine $x$.

\textbf{Resolução:}

A moda é o valor que mais se repete.
No conjunto, $x$ aparece duas vezes, enquanto os demais aparecem uma única vez.

Logo, $x$ pode assumir qualquer valor diferente de 4, 6 e 8.

\textbf{Resposta:} $x \neq 4, 6, 8$.

% QUESTÃO 6
\item A amplitude do conjunto $(x, 4, 11)$ é 8. Determine $x$. Considere que o conjunto está em ordem crescente.

\textbf{Resolução:}

Amplitude:
\[
A = x_{\text{máx}} - x_{\text{mín}}
\]

Para obter amplitude 8:
\[
11 - x = 8
\Rightarrow x = 3 \]

\textbf{Resposta:} $x = 3$ .

% QUESTÃO 7
\item O conjunto $(x, 4, 6, 10)$ tem média igual a 6. Determine $x$.

\textbf{Resolução:}

\[
\frac{x + 4 + 6 + 10}{4} = 6
\Rightarrow \frac{x + 20}{4} = 6
\Rightarrow x + 20 = 24
\Rightarrow x = 4
\]

\textbf{Resposta:} $x = 4$.

% QUESTÃO 8
\item O desvio padrão do conjunto $(5, 5, x)$ é zero. Determine $x$.

\textbf{Resolução:}

Desvio padrão zero indica que todos os valores são iguais.

Logo:
\[
x = 5
\]

\textbf{Resposta:} $x = 5$.

% QUESTÃO 9
\item O conjunto $(x, x+2, x+4)$ possui média igual a 12. Determine $x$.

\textbf{Resolução:}

\[
\frac{x + (x+2) + (x+4)}{3} = 12
\Rightarrow \frac{3x + 6}{3} = 12
\Rightarrow 3x + 6 = 36
\Rightarrow 3x = 30
\Rightarrow x = 10
\]

\textbf{Resposta:} $x = 10$.

% QUESTÃO 10
\item A variância do conjunto $(6, x, 10)$ é 8. Determine $x$.

\textbf{Resolução:}

Média:
\[
\bar{x} = \frac{6 + x + 10}{3} = \frac{x + 16}{3}
\]

Variância:
\[
\frac{(6 - \bar{x})^2 + (x - \bar{x})^2 + (10 - \bar{x})^2}{3} = 8
\]

Resolvendo a equação, obtemos:
\[
x = 4 \text{ ou } x = 12
\]

\textbf{Resposta:} $x = 4$ ou $x = 12$.

\end{enumerate}

\end{multicols}

\end{document}
